\section{硬件与平台：NVIDIA 的护城河与下一轮分工}

\subsection{CUDA 生态为何仍是核心优势}

访谈对 NVIDIA 的判断并不神秘：真正壁垒是两十年积累的 CUDA 生态与开发者工具，而不只是单代芯片参数。即便 TPU、Trainium 与各类 ASIC 快速推进，通用性和生态惯性仍让 NVIDIA 在高变化期保持优势。

\begin{knowledgebox}{训练与推理分离正在加速}
随着推理需求爆发，硬件分工会越来越明确：训练追求带宽与并行效率，推理追求功耗与成本效率。这个趋势会推动更多专用推理芯片出现，也会倒逼软件栈做更细的调度与路由优化。
\end{knowledgebox}

\begin{importantbox}{高增速时期的平台优势逻辑}
只要模型迭代速度快，最有价值的平台通常是“最灵活、最能快速支持新工作负载”的平台，而不一定是单项峰值最佳的平台。NVIDIA 当前仍占据这个位置，但压力在上升。
\end{importantbox}

\begin{warningbox}{“硬件替代”通常慢于媒体叙事}
即便新芯片在某个指标领先，完整迁移还要面对编译链、算子适配、工程团队学习曲线和线上风险控制。平台切换是多年工程，不是季度新闻。
\end{warningbox}

\subsection{本章小结}

硬件竞争会更激烈，但软件生态与迁移成本决定了领先地位不会瞬间翻转。看硬件格局必须同时看工具链与开发者迁移路径。

